\documentclass[11pt]{article}
\usepackage[margin=1in]{geometry}
\usepackage{amsmath,amssymb,amsthm,hyperref}
\newtheorem{theorem}{Theorem}
\title{Linear growth under an unbounded additive driver}
\author{ziangni-sys}
\date{Conjecture 00000008480}
\begin{document}
\sloppy
\maketitle
\noindent The source defines nonautonomous optimization through the averages
$n^{-1}\sum_i(f(T^i x)+h_i)$ and asserts sublinear growth of their supremum
under unbounded driving. Here the supremum is over the initial state $x$
for the displayed finite averaging window. We disprove this assertion using
the unbounded driver $h_i=i$.

\begin{theorem}
There is a compact probability-preserving ergodic system with a continuous
observable and an unbounded nonnegative driver such that the optimized
averages grow linearly, rather than sublinearly.
\end{theorem}
\begin{proof}
Let $\Omega=\{*\}$ with its usual topology and full sigma-algebra, let
$\mu=\delta_*$, and let $T$ be the identity. This space is compact, $T$ is
continuous and preserves probability, and it is ergodic: its only measurable
sets are $\varnothing$ and $\Omega$. Set $f(*)=0$, a continuous observable,
and $h_i=i$ for every integer $i\geq0$. This driver is nonnegative and
unbounded, since for every $B\in\mathbb R$ there is an integer $i>B$.

For a fixed window origin $k\geq0$ and $n\geq1$, define
\[
 A_{k,n}(x)=\frac1n\sum_{i=0}^{n-1}
       \bigl(f(T^i x)+h_{k+i}\bigr),\qquad
 M_{k,n}=\sup_{x\in\Omega} A_{k,n}(x).
\]
The source's displayed average is the case $k=0$. Since $\Omega$ has one
element, the set being optimized is the singleton $\{A_{k,n}(*)\}$;
its supremum is attained and equals that number. The arithmetic sum gives
\[
 M_{k,n}=\frac1n\sum_{i=0}^{n-1}(k+i)
 =\frac{nk+n(n-1)/2}{n}=k+\frac{n-1}{2}.
\]
Thus, for every fixed $k$,
\[
 \frac{M_{k,n}}{n}
 =\frac12+\left(k-\frac12\right)\frac1n
 \longrightarrow\frac12\ne0.
\]
In particular, $M_{0,n}=(n-1)/2$ is not $o(n)$. The claimed universal
sublinear growth under unbounded driving therefore fails.
\end{proof}

\section*{Scope}
Both the English and Chinese versions state the unbounded-driver clause
without a growth or cancellation restriction on the driver. Our driver is
as simple as a nonnegative linear sequence; the autonomous observable is
bounded and identically zero. The example satisfies the usual compactness,
continuity, probability preservation and ergodicity conditions as well.
The proof refutes the sublinear clause, so no claim about the other bounded
or random-driver clauses is required.

\newpage
\section*{Which supremum is computed}
The optimized set here consists of values over initial states for a given
finite $n$-term average, exactly as specified by the displayed expression in
the source. The extra parameter $k$ shows that the counterexample works at
every fixed starting time. It is held fixed while $n$ tends to infinity.
We do not identify this finite-window supremum with a different operation
that also takes a supremum over all starting times. Such an additional
operation is not needed for the counterexample.

\section*{Formal correspondence}
The Lean project defines the actual one-point Dirac probability measure and
identity map. It proves \texttt{CompactSpace}, continuity of both $T$ and
$f$, and Mathlib's actual \texttt{Ergodic} predicate, which includes measure
preservation. The theorem \texttt{driver\_unbounded} states
$\forall B\in\mathbb R,\ \exists i\in\mathbb N,\ B<h_i$.

The definition \texttt{average} uses an actual finite sum over
\texttt{Finset.range n}, actual iterates \texttt{T\^{}[i]}, the observable,
and the shifted driver. The definition \texttt{optimal} is the actual real
supremum \texttt{sSup (Set.range (average start n))}. The complete image set
is proved to be a singleton in \texttt{average\_range};
\texttt{optimal\_attained} then uses the genuine singleton supremum theorem.
Thus no scalar formula for the supremum is assumed.

The arithmetic identity is proved by induction on the natural index in
\texttt{sum\_formula}. The theorem \texttt{shifted\_sum} accounts for every
term of the shifted window, and \texttt{optimal\_formula} establishes
\[
 \forall k,n\in\mathbb N,\quad n\geq1\ \Longrightarrow\
 M_{k,n}=k+(n-1)/2.
\]
The real reciprocal limit gives \texttt{ratio\_tendsto\_half}.
Finally, \texttt{not\_sublinear} proves, for every fixed natural $k$, the
negation of Mathlib's actual \texttt{Asymptotics.IsLittleO} predicate for
\texttt{optimal k} relative to $n\mapsto n$ at infinity. If that predicate
held, the ratio would tend to zero, contradicting uniqueness of its proved
limit $1/2$.

The index zero is harmless: Lean uses totalized division there, whereas
every displayed finite-window formula requires $n\geq1$. All convergence
claims are at infinity, and are unaffected by that single initial index.

\section*{Reproduction}
The project pins Lean 4.19.0 and Mathlib commit
\begin{center}\small\texttt{c44e0c8ee63ca166450922a373c7409c5d26b00b}.\end{center}
From \texttt{lean/}, run \texttt{lake update}, \texttt{lake exe cache get},
and \texttt{lake build}. Check the complete source with
\begin{center}\small\texttt{lake env lean -DwarningAsError=true Main.lean}.\end{center}
Six printed dependency audits cover the system, driver, actual supremum,
formula, limit and final negation. Only standard logical axioms are used.
No admitted statements, extra axioms or auxiliary numerical computations
are needed.
\end{document}
